% SJTU thesis style reconstructed from the Word template comments.
\usepackage[a4paper,top=3.5cm,bottom=4cm,left=3cm,right=2.5cm,headheight=15pt,headsep=0.75cm,footskip=1cm]{geometry}
\usepackage{fontspec}
\usepackage{xeCJK}
\usepackage{amsmath,amssymb,bm}
\usepackage{graphicx}
\usepackage{booktabs,longtable,array,tabularx}
\usepackage{caption}
\usepackage{fancyhdr}
\usepackage{titlesec}
\usepackage{titletoc}
\usepackage{tocloft}
\usepackage{setspace}
\usepackage{indentfirst}
\usepackage{enumitem}
\usepackage{hyperref}
\usepackage{etoolbox}
\usepackage{float}

\IfFontExistsTF{Times New Roman}{\setmainfont{Times New Roman}}{\setmainfont{TeX Gyre Termes}}
\IfFontExistsTF{SimSun}{\setCJKmainfont[AutoFakeBold=true]{SimSun}}{\setCJKmainfont[AutoFakeBold=true]{FandolSong-Regular}}
\IfFontExistsTF{SimHei}{\setCJKfamilyfont{hei}[AutoFakeBold=true]{SimHei}}{\setCJKfamilyfont{hei}[AutoFakeBold=true]{FandolHei-Regular}}
\IfFontExistsTF{SimSun}{\setCJKfamilyfont{song}[AutoFakeBold=true]{SimSun}}{\setCJKfamilyfont{song}[AutoFakeBold=true]{FandolSong-Regular}}

\hypersetup{colorlinks=true,linkcolor=black,citecolor=black,urlcolor=black}
\setcounter{secnumdepth}{3}
\setcounter{tocdepth}{3}
\numberwithin{equation}{chapter}
\numberwithin{figure}{chapter}
\numberwithin{table}{chapter}
\renewcommand{\thefigure}{\thechapter.\arabic{figure}}
\renewcommand{\thetable}{\thechapter.\arabic{table}}
\renewcommand{\theequation}{\thechapter.\arabic{equation}}

\ctexset{
  chapter={name={第,章},number=\arabic{chapter},format={\centering\heiti\bfseries\zihao{3}},beforeskip=24bp,afterskip=18bp,fixskip=true},
  section={format={\raggedright\heiti\bfseries\zihao{4}},beforeskip=24bp,afterskip=6bp,fixskip=true},
  subsection={format={\raggedright\heiti\bfseries\zihao{-4}},beforeskip=12bp,afterskip=6bp,fixskip=true},
  subsubsection={format={\raggedright\songti\zihao{-4}},beforeskip=6bp,afterskip=6bp,fixskip=true}
}

\zihao{-4}\linespread{1.54}\selectfont
\setlength{\parindent}{2em}
\setlength{\parskip}{0pt}
\AtBeginDocument{\zihao{-4}\linespread{1.54}\selectfont}

\pagestyle{fancy}
\fancyhf{}
\fancyhead[L]{\zihao{5}上海交通大学硕士学位论文}
\fancyhead[R]{\zihao{5}\leftmark}
\fancyfoot[C]{\zihao{5}\thepage}
\renewcommand{\headrulewidth}{0.4pt}
\fancypagestyle{plain}{\fancyhf{}\fancyfoot[C]{\zihao{5}\thepage}\renewcommand{\headrulewidth}{0pt}}
\newcommand{\frontplain}{\pagestyle{plain}}
\newcommand{\mainstyle}{\pagestyle{fancy}}

\renewcommand{\contentsname}{目\quad 录}
\titlecontents{chapter}[0em]{\songti\bfseries\zihao{-4}}{\thecontentslabel\quad}{ }{\titlerule*[0.6pc]{.}\contentspage}
\titlecontents{section}[2em]{\songti\zihao{-4}}{\thecontentslabel\quad}{ }{\titlerule*[0.6pc]{.}\contentspage}
\titlecontents{subsection}[4em]{\songti\zihao{-4}}{\thecontentslabel\quad}{ }{\titlerule*[0.6pc]{.}\contentspage}

\captionsetup[figure]{font={small},labelfont={bf},labelsep=quad,justification=centering,position=bottom}
\captionsetup[table]{font={small,bf},labelsep=quad,justification=centering,position=top,skip=6bp}
\renewcommand{\figurename}{图}
\renewcommand{\tablename}{表}

\newcommand{\thesistitlezh}{基于极化带效应水下目标探测感知系统的艇体目标探测研究}
\newcommand{\thesistitleen}{Research on Hull Target Detection Based on an Underwater Target Detection and Sensing System Using Polarization Band Effect}
\newcommand{\thesisauthor}{程家诺}
\newcommand{\thesisstudentid}{124039910036}
\newcommand{\thesissupervisor}{侯中宇}
\newcommand{\thesisschool}{集成电路学院}
\newcommand{\thesismajor}{集成电路工程}
\newcommand{\thesisdegree}{工学硕士}
\newcommand{\thesisdate}{2025年11月2日}

\newcommand{\sjtuFrontTitle}[1]{{\centering\heiti\bfseries\zihao{3}#1\par\vspace{1.5em}}}
\newcommand{\keywords}[1]{\vspace{1em}\noindent\textbf{关键词：}#1\par}
\newcommand{\enkeywords}[1]{\vspace{1em}\noindent\textbf{Key words: }#1\par}
\newenvironment{refsectionmanual}{\chapter*{参考文献}\addcontentsline{toc}{chapter}{参考文献}\zihao{5}\linespread{1.23}\selectfont\begin{enumerate}[label={[\arabic*]},leftmargin=2.6em,itemsep=3bp,topsep=0pt]}{\end{enumerate}}
